\documentclass{slixte}



%%%%%%%%%%%%%%%%%%%%%%%%%%%%%%
%%% Customizing the colors %%%
%%%%%%%%%%%%%%%%%%%%%%%%%%%%%%

% \definecolor{page}{RGB}{255, 255, 255}
% \definecolor{text}{RGB}{0, 0, 0}
% \definecolor{first}{RGB}{96, 159, 103}
% \definecolor{second}{RGB}{103, 96, 159}
% \definecolor{third}{RGB}{159, 103, 96}



%%%%%%%%%%%%%%%%%%%%%%%%%%%%%%%%%%
%%% Customizing the background %%%
%%%%%%%%%%%%%%%%%%%%%%%%%%%%%%%%%%

% \showbackground
% \hidebackground

% \def\backgroundanchor{...}

% \newlength{\backgroundx}
% \setlength{\backgroundx}{...}
% \newlength{\backgroundy}
% \setlength{\backgroundy}{...}
% \newlength{\backgroundwidth}
% \setlength{\backgroundwidth}{...}
% \newlength{\backgroundheight}
% \setlength{\backgroundheight}{...}

%%%%%%%%%%%%%%%%%%%%%%%%%%%%
%%% Customizing the logo %%%
%%%%%%%%%%%%%%%%%%%%%%%%%%%%

% \showlogo
% \hidelogo

% \def\logoanchor{...}

% \newlength{\logox}
% \setlength{\logox}{...}
% \newlength{\logoy}
% \setlength{\logoy}{...}
% \newlength{\logowidth}
% \setlength{\logowidth}{...}
% \newlength{\logoheight}
% \setlength{\logoheight}{...}

%%%%%%%%%%%%%%%%%%%%%%%%%%%%%
%%% Customizing the title %%%
%%%%%%%%%%%%%%%%%%%%%%%%%%%%%

% \setfontsizetitle{...}

% \def\titlealign{...}
% \def\titleanchor{...}

% \newlength{\titlex}
% \setlength{\titlex}{...}
% \newlength{\titley}
% \setlength{\titley}{...}
% \newlength{\titlewidth}
% \setlength{\titlewidth}{...}

% \colorlet{title}{...}

%%%%%%%%%%%%%%%%%%%%%%%%%%%%%%%%
%%% Customizing the subtitle %%%
%%%%%%%%%%%%%%%%%%%%%%%%%%%%%%%%

% \setfontsizesubtitle{...}

% \def\subtitlealign{...}
% \def\subtitleanchor{...}

% \newlength{\subtitlex}
% \setlength{\subtitlex}{...}
% \newlength{\subtitley}
% \setlength{\subtitley}{...}
% \newlength{\subtitlewidth}
% \setlength{\subtitlewidth}{...}

% \colorlet{subtitle}{...}

%%%%%%%%%%%%%%%%%%%%%%%%%%%%%%
%%% Customizing the author %%%
%%%%%%%%%%%%%%%%%%%%%%%%%%%%%%

% \setfontsizeauthor{...}

% \def\authoralign{...}
% \def\authoranchor{...}

% \newlength{\authorx}
% \setlength{\authorx}{...}
% \newlength{\authory}
% \setlength{\authory}{...}
% \newlength{\authorwidth}
% \setlength{\authorwidth}{...}

% \colorlet{author}{...}

%%%%%%%%%%%%%%%%%%%%%%%%%%%%
%%% Customizing the info %%%
%%%%%%%%%%%%%%%%%%%%%%%%%%%%

% \setfontsizeinfo{...}

% \def\infoalign{...}
% \def\infoanchor{...}

% \newlength{\infox}
% \setlength{\infox}{...}
% \newlength{\infoy}
% \setlength{\infoy}{...}
% \newlength{\infowidth}
% \setlength{\infowidth}{...}

% \colorlet{info}{...}



%%%%%%%%%%%%%%%%%%%%%%%%%%%%%%
%%% Customizing the banner %%%
%%%%%%%%%%%%%%%%%%%%%%%%%%%%%%

% \showbanner
% \hidebanner
% \showtopbanner
% \hidetopbanner
% \showbottombanner
% \hidebottombanner
% \showprogress
% \hideprogress

% \setfontsizetopbanner{...}
% \setfontsizebottombanner{...}

% \def\banneropacity{...}
% \def\topbanneropacity{...}
% \def\bottombanneropacity{...}
% \def\progressopacity{...}

% \newlength{\bannerheight}
% \setlength{\bannerheight}{...}
% \newlength{\topbannerheight}
% \setlength{\topbannerheight}{...}
% \newlength{\bottombannerheight}
% \setlength{\bottombannerheight}{...}
% \newlength{\progressheight}
% \setlength{\progressheight}{...}

% \colorlet{banner}{...}
% \colorlet{topbanner}{...}
% \colorlet{bottombanner}{...}
% \colorlet{progress}{...}

%%%%%%%%%%%%%%%%%%%%%%%%%%%%%%%%%%%
%%% Customizing the banner text %%%
%%%%%%%%%%%%%%%%%%%%%%%%%%%%%%%%%%%

% \showbannertext
% \hidebannertext
% \showbannerheader
% \hidebannerheader
% \showbannertitle
% \hidebannertitle
% \showbannerauthor
% \hidebannerauthor
% \showbannersection
% \hidebannersection
% \showbannernumber
% \hidebannernumber

% \def\bannerheaderalign{...}
% \def\bannertitlealign{...}
% \def\bannerauthoralign{...}
% \def\bannersectionalign{...}
% \def\bannernumberalign{...}

% \colorlet{bannertext}{...}
% \colorlet{bannerheader}{...}
% \colorlet{bannertitle}{...}
% \colorlet{bannerauthor}{...}
% \colorlet{bannersection}{...}
% \colorlet{bannernumber}{...}



%%%%%%%%%%%%%%%%%%%%%%%%%%%%%%%%%%%%%%%%%
%%% Customizing the table of contents %%%
%%%%%%%%%%%%%%%%%%%%%%%%%%%%%%%%%%%%%%%%%

% \showtoc
% \hidetoc

% \def\tocopacity{...}

% \newlength{\tocy}
% \setlength{\tocy}{...}
% \newlength{\tocheight}
% \setlength{\tocheight}{...}

%%%%%%%%%%%%%%%%%%%%%%%%%%%%%%%
%%% Customizing the section %%%
%%%%%%%%%%%%%%%%%%%%%%%%%%%%%%%

% \showsection
% \hidesection

% \setfontsizesection{...}

% \def\sectionalign{...}
% \def\sectionanchor{...}

% \newlength{\sectionx}
% \setlength{\sectionx}{...}

% \colorlet{section}{...}
% \colorlet{tosection}{...}

%%%%%%%%%%%%%%%%%%%%%%%%%%%%%%
%%% Customizing the symbol %%%
%%%%%%%%%%%%%%%%%%%%%%%%%%%%%%

% \showsymbol
% \hidesymbol

% \setfontsizesymbol{...}

% \def\symbolalign{...}
% \def\symbolanchor{...}

% \newlength{\symbolx}
% \setlength{\symbolx}{...}

% \colorlet{symbol}{...}



%%%%%%%%%%%%%%%%%%%%%%%%%%%%%%%%%%%%%%%
%%% Customizing the general content %%%
%%%%%%%%%%%%%%%%%%%%%%%%%%%%%%%%%%%%%%%

% \setfontsizecontent{...}
% \setfontsizeref{...}

% \def\contentalign{...}
% \def\contentopacity{...}

% \setlength{\contentmargin}{...}

% \colorlet{item}{...}
% \colorlet{framing}{...}
% \colorlet{ref}{...}
% \colorlet{final}{...}

% \spliton
% \splitoff



%%%%%%%%%%%%%%%%%%%%%%%%%%%%
%%% Fundamental commands %%%
%%%%%%%%%%%%%%%%%%%%%%%%%%%%

\background{path/to/file.png}
\hidebackground % comment out to show the background

\logo{path/to/file.png}
\hidelogo % comment out to show the logo

\title{Title of the slides}
\subtitle{The subtitle of the slides}
\author{Author Name}
\info{Extra information}

\addsection[\textbf{1.}]{First section: basic commands}
\addsection[\textbf{2.}]{Second section: framed environments}
\addsection[\textbf{3.}]{Third section: references and links}
\addsection[\textbf{4.}]{Fourth section: splitting the slides}

\addref{1}{Author(s)}{Year}{Title}{Journal}
\addref{2}{Author(s)}{Year}{Title}{Journal}
\addref{3}{Author(s)}{Year}{Title}{Journal}
\addref{4}{Author(s)}{Year}{Title}{Journal}
\addref{5}{Author(s)}{Year}{Title}{Journal}



\begin{document}

\titleslide

\toc%[Title of table of contents]

\section

\begin{tikzslide}[A tikzslide]
    \draw (\paperwidth/2, \paperheight/2) circle (5cm);
    \content{
        Some text within the tikzslide.
    }
\end{tikzslide}

\begin{slide}[A standard  text slide]
    You can easily create a list.
    \begin{itemize}
        \fillitem A first element.
        \fillitem[second] A second element.
        \begin{itemize}
            \emptyitem A first sub-element.
            \emptyitem[second] A second sub-element.
        \end{itemize}
        \citem[third]{\textbf{?}} A third and different element.
    \end{itemize}
\end{slide}

\begin{slide}[A slide with useful environments]
    \begin{alert}%[...]
        An important information in the alert environment.
    \end{alert}

    \spacing

    \begin{remark}%[...]
        Something worth knowing in the remark environment.
    \end{remark}

    \spacing

    \begin{pointer}%[...]
        A highlighted information in the pointer environment.
    \end{pointer}

    \spacing

    \begin{question}%[...]
        A question in the question environment.
    \end{question}
\end{slide}

\section

\begin{slide}[A slide with basic framed environments]
    \begin{definition}%[...]
        An important definition
    \end{definition}
    \begin{theorem}%[...]
        The statement of a related theorem.
    \end{theorem}
\end{slide}

\begin{slide}[A slide with a more advanced framed environments]
    \begin{theorem}[Your Initial, 2025]
        The statement of your new theorem.
    \end{theorem}
    \begin{problem}%[...]
        Something still unknown.
    \end{problem}
\end{slide}

\begin{slide}[A slide with the most general framed environment]
    \begin{result}{Proposition}{}
        A proposition related to your presentation.
    \end{result}
    \begin{result}{Corollary}{From previous proposition}
        A corollary based on your previous proposition.
    \end{result}
\end{slide}

\section

\begin{slide}[References and links]
    You can easily create hyperlinks to various parts of the slides.
    \begin{itemize}
        \fillitem Some references, such as the first one\citeref{1}, the second one\citeref{2}, etc.
        \fillitem Your current section, also called \tosection{}.
        \fillitem Other sections: \tosection{1} or \tosection{4}.
        \fillitem It is also worth mentioning other implemented links within the slides.
        \begin{itemize}
            \emptyitem Clicking on the title (usually bottom left) brings back to the title page.
            \emptyitem Clicking on the author (usually bottom right) brings to the final page (if defined).
            \emptyitem Clicking on the section (usually bottom center) brings to the beginning of the section.
            \emptyitem Clicking on the sections in the table of contents brings to their beginnings.
            \emptyitem Clicking on a reference within a slide brings to the reference page (if defined).
            \emptyitem Clicking on a reference within the reference slide brings you to the first use of the reference (if defined).
        \end{itemize}
    \end{itemize}
\end{slide}

\section

\begin{slide}[A linear slide]
    One might want to create a slide where points appear after each other.
    \begin{itemize}
        \break\fillitem This appears first.
        \break\fillitem Then this appears second.
        \break\fillitem And finally this appears at the end.
    \end{itemize}

    \spacing

    It is also worth noting the existence of a default space command to break the page.
    
    \spacing[2]

    This spacing can also be customized.
\end{slide}

\begin{slide}[A breaking slide]
    You can split a single slide to create effects and keep the presentation flowing.
    \begin{itemize}
        \makeon[1]{\fillitem Some text not taking space when not on the slide.}
        \showon[-2]{\fillitem Some text taking space even when not on the slide.}
        \fillitem Some text on all slides to show its interaction with the other parts.
    \end{itemize}
    \showon[3]{}
\end{slide}

\refslide%[Title of reference page]

\finalslide%[Final message]

\end{document}